\begin{thebibliography}{9}
\bibitem{proveit} V. Reshetnikov et al., \emph{ProveIt}, Logic/PrimitiveRecursiveArithmetic. Ordinary iPRA is defined in \texttt{ConciseIPRA.PR}; the integer, order, power and bounded-sum modules are in \texttt{IPRADerivations.Core}. Repository snapshot \texttt{84192061db7eb09e923bd4ded5199f23d30c6e83}. \href{https://github.com/VladimirReshetnikov/ProveIt/tree/84192061db7eb09e923bd4ded5199f23d30c6e83/Logic/PrimitiveRecursiveArithmetic}{Pinned source and verification-status guide}.
\bibitem{busch} P. Busch, Quantum states and generalized observables: a simple proof of Gleason's theorem. \emph{Phys. Rev. Lett.} \textbf{91} (2003), 120403. Classical context for the all-effects additivity assumption in Theorem~\ref{thm:q-effect}. \href{https://doi.org/10.1103/PhysRevLett.91.120403}{doi:10.1103/PhysRevLett.91.120403}.
\bibitem{watrous} J. Watrous, \emph{The Theory of Quantum Information}, Cambridge University Press, 2018, chapters 2--3. Finite channel and discrimination terminology. \href{https://cs.uwaterloo.ca/~watrous/TQI/}{Author's text and errata}.
\bibitem{lindblad} G. Lindblad, On the generators of quantum dynamical semigroups. \emph{Comm. Math. Phys.} \textbf{48} (1976), 119--130. Classical generator theorem underlying the finite-dimensional formulation of Section~\ref{sec:q-open}. \href{https://doi.org/10.1007/BF01608499}{doi:10.1007/BF01608499}.
\bibitem{brattka} V. Brattka and R. Dillhage, Computability of the spectrum of self-adjoint operators. \emph{J. Universal Computer Science} \textbf{11} (2005), 1884--1900. Computability of spectral information for different operator presentations. \href{https://doi.org/10.3217/jucs-011-12-1884}{doi:10.3217/jucs-011-12-1884}.
\bibitem{ye} F. Ye, Toward a constructive theory of unbounded linear operators. \emph{J. Symbolic Logic} \textbf{65} (2000), 357--370. Constructive spectral and Stone theorems in Bishop's setting. \href{https://doi.org/10.2307/2586543}{doi:10.2307/2586543}.
\bibitem{hellman} G. Hellman, Constructive mathematics and quantum mechanics: unbounded operators and the spectral theorem, in \emph{Mathematics and Its Logics}, Cambridge University Press, 2021, pp. 212--236. Background on the constructive representation problem. \href{https://doi.org/10.1017/9781108657419.014}{doi:10.1017/9781108657419.014}.
\end{thebibliography}
